\chapter{Chapter 12 Refresher: Modern LLM Access --- APIs, SDKs, Agents, and the Model Context Protocol --- Latency, Queues, and Reliability Budgets}

\section*{Setting the Scene}
APIs expose stochastic models inside stochastic infrastructure: arrivals, service times, retries, and pricing multiply Preface Steps~5--7 expectations into tail experiences users actually feel.

\paragraph{Step A --- Utilisation as latent probability.}
\textbf{Problem.} Saturated GPUs inflate latency although average load looks fine.
\textbf{Insight.} \(\rho=\lambda/\mu\) nearing 1 blows variance (Preface Step~6)---queue waits explode superlinearly.
\textbf{Lineage.} Kendall queue notation; Kleinrock textbooks.
\textbf{Mental model.} Treat \(\rho\) like criticality in statistical mechanics---small deltas matter near phase transitions.

\paragraph{Step B --- Retries compound load.}
\textbf{Problem.} Clients retry on flakes.
\textbf{Insight.} Geometric retry counts lift effective \(\lambda\)---same algebra as Preface Step~4 inverted conditioning stories but operational.

\paragraph{Step C --- Guards cap rare disasters.}
\textbf{Problem.} Agents might loop unbounded.
\textbf{Insight.} Hard budgets reinstate finite \(\Omega\) partitions compatible with Step~1 measurability mindset.

Probability-minded reliability engineers read LM dashboards the same way traders read risk tapes---moments plus tails.

\section*{Notation Introduced in This Chapter}
\begin{itemize}[leftmargin=1.5em]
  \item $\lambda$: request arrival rate.
  \item $\mu$: service rate.
  \item $\rho=\lambda/\mu$: utilization.
  \item $\mathbb{E}[W]$: expected waiting time.
  \item $p_f$: per-attempt failure probability.
  \item $\mathbb{E}[K]$: expected number of attempts.
\end{itemize}

\section*{What You Need To Recall Precisely}
\begin{itemize}[leftmargin=1.5em]
  \item \textbf{Latency decomposition:} prefill latency plus per-output-token decode latency.
  \item \textbf{Queueing baseline (M/M/1).}
  Assume Poisson arrivals rate \(\lambda\), exponential service rate \(\mu\) per busy server, single server, infinite buffer (steady-state regime \(\lambda<\mu\)). Expected waiting time in queue satisfies
  \[
    \mathbb{E}[W]=\frac{1}{\mu-\lambda},
  \]
  blowing up as \(\rho=\lambda/\mu\uparrow 1$: tiny overload deltas dominate latency tails---engineering intuition beats pretending arrivals are perfectly deterministic.

  \item \textbf{Retry expectation.}
  If each attempt fails independently with probability \(p_f\) and retries repeat until first success (no cap), successes follow geometric distribution with success probability \(1-p_f\). Expected attempts \(\mathbb{E}[K]=\sum_{n\ge1} n\, p_f^{n-1}(1-p_f)=1/(1-p_f)\). Even modest \(p_f\) inflates load multipliers \(\mathbb{E}[K]\) nonlinearly.
  \item \textbf{Token cost model:} base cost is linear in tokens, expected cost depends on retries and aborts.
  \item \textbf{Termination guards:} bounded steps/tokens/cost enforce finite agent behavior.
\end{itemize}

\section*{How These Results Build on Each Other}
\begin{enumerate}[leftmargin=1.5em]
  \item Estimate traffic intensity $(\lambda/\mu)$.
  \item Predict queue delay and saturation risk.
  \item Incorporate retries into effective load and expected cost.
  \item Enforce runtime guards for bounded execution.
\end{enumerate}

\section*{Reminder Glossary (Term by Term)}
\begin{itemize}[leftmargin=1.5em]
  \item \textbf{Arrival rate $\lambda$:} requests per unit time.
  \item \textbf{Service rate $\mu$:} processing capacity per unit time.
  \item \textbf{Utilization $\rho$:} $\lambda/\mu$.
  \item \textbf{Retry multiplier:} expected attempts per successful request.
  \item \textbf{Budget guard:} hard bound on cost/tokens/steps.
\end{itemize}

\section*{Worked Mini Example (Retries)}
Let per-attempt failure probability be $p_f=0.2$ and assume independent retries with no cap.
\[
\mathbb{E}[K]=\frac{1}{1-p_f}=\frac{1}{0.8}=1.25.
\]
If average successful request cost is \$0.04, expected cost with retries is roughly
\[
0.04\times 1.25 = \$0.05.
\]

\section*{Worked Example (Queue and Retry Coupling)}
If base arrival rate is 80 req/s and retries add 25\% expected attempts, effective arrival becomes:
\[
\lambda_{\text{eff}} = 80 \times 1.25 = 100\ \text{req/s}.
\]
If service rate is $\mu=110$ req/s, utilization rises from $0.73$ to $0.91$.  
Interpretation: retry policy can push system into tail-latency regime even when baseline looked safe.

\section*{Intuition Checkpoints}
\begin{itemize}[leftmargin=1.5em]
  \item Systems can meet average SLOs while failing badly in tails.
  \item Queueing approximations are useful planning tools, not exact traffic simulators.
  \item Reliability controls (backoff, circuit breaking) alter both latency and cost.
\end{itemize}

\section*{Further Refresh}
\begin{itemize}[leftmargin=1.5em]
  \item Systems scaling and bottlenecks: see Chapter 7 refresher.
  \item Evaluation and trust boundaries: see Chapter 10 refresher.
  \item Risk thresholding and controls: see Chapter 11 refresher.
  \item \textbf{Harchol-Balter} --- queueing and latency scaling.
  \item \textbf{SRE workbooks} --- retries/backoff/overload control design.
\end{itemize}
